```latex
\chapter{Introducción general}
\label{chap:introduction}

\section{Propósito y contexto del proyecto}

Este documento constituye el informe final del sistema de monitoreo espectral \textbf{ANE-HX}, desarrollado como proyecto académico de ingeniería orientado a la documentación técnica, integración, validación y transferencia de conocimiento de una solución de monitoreo del espectro radioeléctrico.

El sistema integra componentes de hardware, adquisición por radio definida por software (SDR), procesamiento en borde, conectividad, plataforma web en la nube, almacenamiento de datos, visualización, análisis espectral, protocolos de prueba y documentación de soporte. Su propósito es consolidar en un único informe la arquitectura del sistema, las decisiones de diseño, los módulos implementados, los problemas encontrados, las pruebas realizadas y las recomendaciones para continuidad o mantenimiento.

El informe se organiza como memoria técnica del sistema y como instrumento de transferencia de conocimiento. Por tanto, no se limita a describir el resultado final, sino que también registra criterios de diseño, limitaciones, ajustes realizados, evidencias de prueba, anexos técnicos y documentos externos de apoyo.

\section{Problema o necesidad principal}

El monitoreo del espectro radioeléctrico requiere capturar señales en campo, procesarlas, extraer parámetros técnicos, transportar información hacia una plataforma de análisis, almacenar resultados, visualizar mediciones y contrastar hallazgos contra criterios técnicos o regulatorios.

Una solución de este tipo debe articular varios subsistemas:

\begin{itemize}[leftmargin=*]
    \item Nodo físico de sensado.
    \item Antenas y cadena de recepción RF.
    \item Equipo SDR para adquisición de señales.
    \item Software de adquisición y preprocesamiento.
    \item Conectividad entre el nodo de borde y la nube.
    \item Plataforma web para operación, administración y visualización.
    \item Base de datos para almacenamiento de mediciones.
    \item Módulos de análisis espectral y localización.
    \item Procedimientos de instalación, prueba, validación y mantenimiento.
\end{itemize}

El proyecto ANE-HX aborda esta necesidad mediante una arquitectura modular que separa el sensado en borde, la adquisición SDR, la transmisión de datos, la persistencia en la nube, el análisis espectral y la interacción con usuarios operadores o administradores.

\section{Alcance del documento}

El presente informe cubre el diseño, implementación, integración, pruebas, documentación de soporte y aprendizajes obtenidos durante el desarrollo del sistema. El documento no pretende reemplazar los repositorios de código fuente, los manuales externos ni los archivos de soporte almacenados fuera del repositorio \LaTeX. En cambio, los referencia y los organiza para facilitar su consulta.

El alcance documental incluye:

\begin{itemize}[leftmargin=*]
    \item Descripción general del sistema y su arquitectura.
    \item Requerimientos, restricciones y criterios de uso.
    \item Diseño físico y electrónico del nodo de sensado.
    \item Software de adquisición SDR y procesamiento en borde.
    \item Plataforma web, servicios cloud, backend, frontend y base de datos.
    \item Análisis espectral, localización y procesamiento de capturas.
    \item Protocolos de prueba y criterios de aceptación.
    \item Integración general del sistema.
    \item Conclusiones, recomendaciones y referencias.
    \item Anexos de apoyo, diagramas, manuales, documentos externos y resultados de prueba.
\end{itemize}

La tabla de contenido, la lista de figuras y la lista de tablas se generan automáticamente durante la compilación del documento. Por tanto, la estructura final del informe debe consultarse principalmente desde dichos índices, ya que se actualizan de acuerdo con los capítulos y anexos efectivamente incluidos desde el archivo \texttt{main.tex}.

\section{Estructura general del informe}

El documento se organiza en capítulos principales y anexos. Cada capítulo se mantiene en un archivo independiente dentro de la carpeta \texttt{section/}, y el archivo \texttt{main.tex} actúa como punto de entrada para compilar el informe completo.

La estructura efectiva del informe no se enumera manualmente en esta sección, porque puede cambiar durante el proceso de edición. Los capítulos y anexos incluidos se determinan directamente desde el archivo \texttt{main.tex}. Por tanto, la referencia oficial a la estructura vigente del documento es la tabla de contenido generada automáticamente al inicio del PDF.

Esta decisión evita inconsistencias cuando se agregan, eliminan, renombran o comentan capítulos que todavía no están completamente desarrollados.

De forma general, el informe contempla los siguientes bloques documentales:

\begin{itemize}[leftmargin=*]
    \item Introducción, contexto, alcance y convenciones del documento.
    \item Requerimientos, restricciones y criterios de uso del sistema.
    \item Diseño general y arquitectura del sistema.
    \item Hardware del nodo de sensado y cadena RF.
    \item Software del sensor, adquisición SDR y procesamiento en borde.
    \item Plataforma web, servicios cloud, backend, frontend y base de datos.
    \item Análisis espectral, localización y procesamiento de capturas.
    \item Protocolos de prueba, validación y criterios de aceptación.
    \item Integración general del sistema.
    \item Conclusiones, recomendaciones y referencias.
    \item Anexos técnicos, diagramas, manuales externos y resultados de pruebas.
\end{itemize}

Esta organización permite que cada autor trabaje sobre su archivo correspondiente sin modificar directamente el cuerpo completo del documento.

Los anexos complementan el cuerpo principal del informe:

\begin{itemize}[leftmargin=*]
    \item \hyperref[chap:ia-collaboration]{Anexo~\ref*{chap:ia-collaboration}: Técnica de desarrollo con agentes de inteligencia artificial}.
    \item \hyperref[chap:additional-diagrams]{Anexo~\ref*{chap:additional-diagrams}: Diagramas adicionales}.
    \item \hyperref[chap:manuals-guides]{Anexo~\ref*{chap:manuals-guides}: Manuales y documentos de apoyo}.
    \item \hyperref[chap:test-results]{Anexo~\ref*{chap:test-results}: Resultados de pruebas}.
\end{itemize}

Esta organización permite que cada autor trabaje sobre su archivo correspondiente sin modificar directamente el cuerpo completo del documento.

\section{Resumen de la arquitectura y módulos principales}

El sistema ANE-HX se organiza en tres dominios funcionales:

\begin{enumerate}[leftmargin=*]
    \item \textbf{Borde:} corresponde al nodo físico de sensado. Incluye el hardware de adquisición, antenas, cadena RF, SDR, tarjeta auxiliar, alimentación, conectividad y software local de control o preprocesamiento.
    
    \item \textbf{Nube:} corresponde a la plataforma central de almacenamiento, visualización, administración y análisis. Incluye servicios de backend, frontend, base de datos, comunicación con sensores, visualización de métricas, gestión de campañas y procesamiento posterior.
    
    \item \textbf{Usuarios:} corresponde a los perfiles que interactúan con la plataforma, principalmente operadores, administradores o usuarios técnicos encargados de consultar mediciones, revisar resultados, gestionar sensores y generar reportes.
\end{enumerate}

A partir de esta división, los capítulos técnicos documentan los subsistemas del proyecto y su relación con el flujo completo de adquisición, transmisión, análisis y visualización de información espectral.

\section{Descripción resumida de los capítulos}

Los capítulos principales documentan los componentes técnicos y metodológicos del sistema. La inclusión final de cada capítulo depende de las instrucciones activas en el archivo \texttt{main.tex}.

\subsection{Requerimientos, restricciones y uso}

Este capítulo consolida los requerimientos funcionales y no funcionales del sistema, las restricciones técnicas, los criterios de uso, las capacidades esperadas y los elementos que permiten contrastar el diseño con el alcance definido para el proyecto.

\subsection{Diseño general del sistema}

Este capítulo describe la arquitectura global del sistema, la separación por subsistemas, el flujo de información, las interfaces principales, las decisiones de diseño y la forma en que los componentes de borde y nube se integran dentro de una solución completa.

\subsection{Hardware del nodo de sensado}

Este capítulo documenta el diseño físico del nodo de sensado, incluyendo componentes electrónicos, alimentación, ruta RF, antenas, conectividad, montaje, integración física y consideraciones de operación en campo.

\subsection{Software del sensor y adquisición SDR}

Este capítulo describe el software ejecutado en el nodo de borde, incluyendo control del SDR, adquisición de muestras, selección de parámetros, preprocesamiento, generación de productos espectrales, almacenamiento temporal y comunicación con la plataforma central.

\subsection{Plataforma web en la nube}

Este capítulo documenta la plataforma web, sus componentes de frontend, backend, base de datos, servicios de comunicación, visualización de datos, administración de sensores, usuarios, campañas y reportes.

\subsection{Análisis espectral y localización}

Este capítulo presenta los métodos de análisis espectral, procesamiento de capturas, extracción de parámetros técnicos, detección de emisiones, localización o asociación geográfica y visualización de resultados analíticos.

\subsection{Protocolos de prueba edge-cloud}

Este capítulo define procedimientos de prueba, criterios de aceptación, validaciones funcionales, pruebas de integración, pruebas de adquisición, pruebas de comunicación y criterios para verificar el comportamiento del sistema.

\subsection{Integración general del sistema}

Este capítulo consolida la integración entre hardware, software de borde, conectividad, plataforma cloud, base de datos, análisis espectral y operación por parte del usuario.

\subsection{Conclusiones generales}

Este capítulo resume los resultados principales del proyecto, los aprendizajes obtenidos, el estado general del sistema, las limitaciones identificadas y los aportes más relevantes del desarrollo.

\subsection{Recomendaciones}

Este capítulo presenta recomendaciones para continuidad, mantenimiento, mejora, pruebas futuras, escalamiento, documentación adicional y posibles líneas de trabajo posteriores.

\subsection{Referencias}

Este capítulo consolida las referencias bibliográficas, normativas, técnicas o documentales utilizadas para respaldar el desarrollo del informe.

\section{Descripción resumida de los anexos}

Los anexos complementan el cuerpo principal del informe. Su inclusión final también depende de las instrucciones activas en el archivo \texttt{main.tex}.

\subsection{Técnica de desarrollo con agentes de inteligencia artificial}

Este anexo documenta la metodología de apoyo con agentes de inteligencia artificial utilizada durante el desarrollo documental, incluyendo criterios de uso, roles, ventajas, limitaciones y aprendizajes.

\subsection{Diagramas adicionales}

Este anexo consolida diagramas complementarios que apoyan la comprensión de la arquitectura, flujos de información, componentes o relaciones entre subsistemas.

\subsection{Manuales y documentos de apoyo}

Este anexo referencia manuales, documentos externos, informes técnicos, videos de ensamble y archivos de soporte disponibles en la carpeta documental del proyecto. Estos documentos son material externo no editable dentro del repositorio \LaTeX.

\subsection{Resultados de pruebas}

Este anexo debe consolidar evidencias, resultados, capturas, salidas, tablas o registros asociados a las pruebas ejecutadas sobre el sistema.

\section{Convenciones del documento}

El informe está redactado en español, compilado con \LaTeX{} y utiliza codificación UTF-8. Cada capítulo debe mantener una estructura clara y homogénea que facilite su revisión, trazabilidad y mantenimiento.

En general, se recomienda que cada capítulo incluya:

\begin{itemize}[leftmargin=*]
    \item Propósito del capítulo.
    \item Contexto dentro del sistema.
    \item Descripción técnica del componente, módulo o proceso.
    \item Entradas, salidas e interfaces.
    \item Decisiones de diseño.
    \item Problemas encontrados y ajustes realizados.
    \item Validación, pruebas o evidencias.
    \item Limitaciones.
    \item Recomendaciones o trabajo futuro.
\end{itemize}

No se debe incluir código fuente extenso dentro del informe. Cuando sea necesario referenciar software, repositorios, scripts o servicios externos, se debe incluir su enlace o indicar explícitamente que el recurso está pendiente de publicación.

\section{Mantenimiento de la estructura documental}

La estructura del informe se controla desde el archivo \texttt{main.tex}. Si un capítulo o anexo todavía no está desarrollado, se recomienda una de las siguientes opciones:

\begin{enumerate}[leftmargin=*]
    \item Mantener el archivo \texttt{.tex} con una nota breve que indique su estado pendiente o parcial.
    \item Retirar temporalmente la línea correspondiente en \texttt{main.tex}.
    \item Usar una instrucción de inclusión que omita archivos inexistentes sin generar capítulos artificiales.
\end{enumerate}

Para evitar inconsistencias, este capítulo no debe listar manualmente el estado de avance de cada archivo. El estado de avance debe mantenerse en el \texttt{README.md}, en una tabla de control documental o en una sección específica de gestión del proyecto.

\section{Recomendación de compilación}

Después de agregar, quitar o renombrar capítulos, anexos, figuras o tablas, se recomienda compilar el documento al menos dos veces. Esto permite actualizar correctamente:

\begin{itemize}[leftmargin=*]
    \item Tabla de contenido.
    \item Lista de figuras.
    \item Lista de tablas.
    \item Numeración de capítulos y anexos.
    \item Referencias cruzadas.
    \item Enlaces internos.
\end{itemize}

Si persisten errores de numeración, nombres automáticos o referencias, se recomienda eliminar los archivos auxiliares de compilación y volver a compilar desde cero.
```
